\section{Cadena de suministro de datos: el proceso industrial que va de la fuente de datos a la capacidad}

Si se concibe los datos como una industria, hace falta una cadena de suministro: recolección, licenciamiento, limpieza, anotación, síntesis, mezcla, entrenamiento, evaluación y retroalimentación. Cada eslabón influye en la capacidad final del modelo. Cuanto más larga es la cadena de suministro de datos, más necesarias se vuelven la estandarización, la auditoría y el control de calidad.

\subsection{Tabla de la cadena de suministro de datos}

\noindent\begin{longtable}{p{0.20\textwidth}p{0.34\textwidth}p{0.37\textwidth}}
\toprule
Eslabón & Tarea & Riesgo de calidad \\
\midrule
Recolección & Obtener datos de páginas web, robots, simulación y comportamiento de usuarios & Sesgo de origen, derechos de autor, privacidad. \\
Limpieza & Eliminar duplicados y ruido, filtrar muestras de baja calidad & Borrar por error casos de cola larga de alto valor. \\
Anotación & Dar estructura a tareas, acciones, estados y retroalimentación & Anotaciones inconsistentes, costo alto. \\
Síntesis & Usar modelos o simulación para cubrir escenarios escasos & Desplazamiento de la distribución, colapso de modos. \\
Mezcla & Definir la proporción de los datos y el orden de entrenamiento & Capacidades desequilibradas, olvido. \\
Evaluación & Revisar las capacidades del modelo y los tipos de falla & Métricas distorsionadas, sobreajuste al benchmark. \\
Retorno & Devolver las fallas y los escenarios nuevos al conjunto de datos & Ciclo cerrado lento, retroalimentación imprecisa. \\
\bottomrule
\end{longtable}

\begin{knowledgebox}{La barrera de entrada de la cadena de suministro de datos}
La barrera de entrada de los datos no consiste solo en «tengo datos», sino también en «sé cómo convertir los datos en capacidad». La limpieza, la proporción de la mezcla, la evaluación y el retorno suelen ser más difíciles de copiar que los datos en bruto.
\end{knowledgebox}

\subsection{Resumen de la sección}

La industrialización del sector de datos implica que entre la fuente de datos y la capacidad existe una cadena de suministro completa. Si cualquier eslabón es débil, la capacidad final del modelo disminuye.
